% ---------------------------------------------------------------------------
% appendix-quickref.tex
% ---------------------------------------------------------------------------
\chapter{One-Page Quick Reference}
\label{app:quickref}

\chapabstract{Every command and file field in this booklet, compressed onto
one page to keep next to the keyboard.}

\section{Commands}

\begin{longtable}{@{}p{0.46\textwidth}p{0.49\textwidth}@{}}
\toprule
\textbf{Command} & \textbf{Purpose}\\
\midrule
\endfirsthead
\toprule
\textbf{Command} & \textbf{Purpose}\\
\midrule
\endhead
\code{fusesoc --cores-root <dir> ...} & Add \code{<dir>} to where \FuseSoC
  looks for \file{.core} files. Repeat once per root you need.\\
\code{fusesoc ... run --target <t> <core>} & Resolve dependencies and run
  target \code{<t>} of \code{<core>}.\\
\code{fusesoc ... run --target <t> <core> --run\_options="..."} & Pass
  arguments straight through to the simulation binary (plusargs).\\
\code{python3 util/vendor.py <descriptor>.vendor.hjson} & Fetch one
  dependency at its pinned revision into \file{vendor/}.\\
\code{make vendor} & Run the vendoring script over every
  \file{*.vendor.hjson} the lab has.\\
\code{make help} & List a lab's own Makefile targets with a one-line
  description each.\\
\code{make <target> --dry-run} & Print the commands a target would run,
  without running them.\\
\code{python3 regtool.py -r -t rtl <regs>.hjson} & Generate
  \code{*\_reg\_pkg.sv} and \code{*\_reg\_top.sv} from a register
  description.\\
\code{python3 regtool.py -D -o <regs>.h <regs>.hjson} & Generate the
  matching C header of offsets, masks and field macros.\\
\code{make regs} & Run both regtool invocations above for a lab that
  defines this target.\\
\bottomrule
\end{longtable}

\section{.core file fields}

\begin{longtable}{@{}p{0.30\textwidth}p{0.65\textwidth}@{}}
\toprule
\textbf{Field} & \textbf{Meaning}\\
\midrule
\endfirsthead
\toprule
\textbf{Field} & \textbf{Meaning}\\
\midrule
\endhead
\code{CAPI=2:} & Required first line; marks the file as Core API v2.\\
\kw{name} & \code{vendor:library:core:version}, the string a dependant
  names in its own \kw{depend} list.\\
\kw{filesets.<name>.files} & Source paths, relative to the \file{.core}
  file, in compile order.\\
\kw{filesets.<name>.file\_type} & \code{systemVerilogSource},
  \code{vhdlSource}, \code{cSource}, \dots\\
\kw{filesets.<name>.depend} & Other cores' names; \FuseSoC resolves them
  recursively.\\
\kw{targets.<name>.filesets} & Which filesets this target pulls in.\\
\kw{targets.<name>.default\_tool} & Which tool block under \kw{tools} to
  use.\\
\kw{targets.<name>.toplevel} & The module to elaborate from.\\
\kw{targets.<name>.tools.<tool>} & Tool-specific options
  (\code{verilator\_options}, \code{vlog\_options}, \dots).\\
\bottomrule
\end{longtable}

\section{.vendor.hjson fields}

\begin{longtable}{@{}p{0.30\textwidth}p{0.65\textwidth}@{}}
\toprule
\textbf{Field} & \textbf{Meaning}\\
\midrule
\endfirsthead
\toprule
\textbf{Field} & \textbf{Meaning}\\
\midrule
\endhead
\kw{name} & A label for the descriptor (not a \FuseSoC core name).\\
\kw{target\_dir} & Where the snapshot is copied to.\\
\kw{upstream.url} & The upstream repository.\\
\kw{upstream.rev} & The pinned tag, branch or commit -- the whole
  reproducibility guarantee.\\
\kw{exclude\_from\_upstream} & Paths dropped from the copy (tests, CI).\\
\bottomrule
\end{longtable}

\section{.hjson register description fields}

\begin{longtable}{@{}p{0.30\textwidth}p{0.65\textwidth}@{}}
\toprule
\textbf{Field} & \textbf{Meaning}\\
\midrule
\endfirsthead
\toprule
\textbf{Field} & \textbf{Meaning}\\
\midrule
\endhead
\kw{name} / \kw{regwidth} & The block's name (prefixes every generated
  symbol) and bus width in bits (normally 32).\\
\kw{registers[].name} / \kw{.desc} & One entry per register; \kw{desc} ends
  up as a comment in the generated header.\\
\kw{swaccess} & Software access: \code{ro}, \code{wo} or \code{rw} --
  what the CPU is allowed to do at this address.\\
\kw{hwaccess} & Hardware access: \code{hro} (hardware reads \kw{.q} only)
  or \code{hwo} (hardware writes \kw{.d} only).\\
\kw{hwqe} & Adds a one-cycle \kw{.qe} write-strobe alongside \kw{.q}, so
  hardware can see \emph{when} software wrote, not just the value.\\
\kw{hwext} & Removes the register's internal storage; \kw{hw2reg} must
  drive the read value combinationally every cycle.\\
\kw{resval} & Reset value, in the generated hardware and in
  documentation.\\
\kw{fields[].bits} / \kw{.name} & Sub-fields of a register, each with its
  own bit range and name (e.g.\ \code{"0"} for \code{DONE}).\\
\bottomrule
\end{longtable}

\section{Writing a new .core file, in five steps}

\begin{enumerate}
  \item Name it: \code{<vendor>:<lab>:<core>:<version>}.
  \item List your sources in a fileset, in compile order.
  \item Declare dependencies by name in \kw{depend}; vendor them first if
        they are not already under \file{vendor/}.
  \item Add one target per tool you run, starting with \code{lint}.
  \item Set \kw{toplevel} to what that target should elaborate.
\end{enumerate}

\section{The three traps to remember}

\begin{enumerate}
  \item A \file{.core} file's filesets do nothing until a target lists
        them; check the target block first when a build seems to be
        missing files.
  \item \kw{upstream.rev: "main"} is fine while developing, wrong once a
        lab ships: pin it to a SHA or tag.
  \item \code{include} of another Makefile must come last if that file
        defines a catch-all rule, or every target after it is invisible.
  \item Generated files (\code{*\_reg\_pkg.sv}, \code{*\_reg\_top.sv},
        \code{*\_regs.h}) are overwritten by \code{make regs}; edit the
        \file{.hjson} source, never the generated files themselves.
\end{enumerate}
